% ANSWER Q6607 — conditional p^{1/6} coincidence theorem (ASCII delivery via message
% fallback; github connector down). Full document as delivered; to be adapted into
% section form + audited before any merge. Received 2026-08-01 evening.
% Key cross-checks vs our banked results:
%   - genus <= 6hk consistent with exact g = 6hk-3h-3k-2 (ledger §112, Q6561)
%   - CENSUS_H = component dichotomy (thm:coinc-components under MON)
%   - honest completion range H <= c p^{1/6}/(2+log p)^{2/3}
%   - explicit: c0=1/4, p>=4096, K=4; corollary sum C_h <= H + 2 sqrt(pH)
\documentclass[11pt]{article}
\usepackage{amsmath,amssymb,amsthm}

\newtheorem{theorem}{Theorem}
\newtheorem{proposition}[theorem]{Proposition}
\newtheorem{lemma}[theorem]{Lemma}
\theoremstyle{remark}
\newtheorem{remark}[theorem]{Remark}

\newcommand{\Fp}{\mathbf F_p}
\newcommand{\Ncoinc}{N_{\mathrm{coinc}}}

\begin{document}

\section*{Conditional $p^{1/6}$ coincidence theorem (ASCII version)}

Let $p \ge 7$ be prime, $M = p-2$, and for $h \ge 1$ define
\[
q_h(X) = \prod_{j=1}^h (X+j), \qquad
\delta_h(X) = \frac{N_h(X)}{q_h(X)^3}.
\]
We work on the nonwrapping windows
\[
I_h = \{1,\dots,M-h\}, \qquad L_h = |I_h| = p-2-h,
\]
and set
\[
S = S_H = \sum_{h=1}^H L_h = pH - \frac{H^2+5H}{2}.
\]

Define
\[
R_{h,k} = \#\{(x,y)\in I_h \times I_k : \delta_h(x) = \delta_k(y)\},
\]
so that
\[
\Ncoinc(H) = \sum_{1 \le h,k \le H} R_{h,k}.
\]

Clearing denominators gives the affine curve
\[
F_{h,k}(X,Y) = N_h(X) q_k(Y)^3 - N_k(Y) q_h(X)^3 = 0,
\]
of bidegree at most $(3h,3k)$.

\medskip

\noindent
\textbf{CENSUS$_H$.}
\begin{itemize}
\item For all $1 \le h < k \le H$, $F_{h,k}$ is absolutely irreducible over $\Fp$.
\item For all $2 \le h \le H$, $F_{h,h}$ has exactly two components:
the diagonal $X=Y$ and an absolutely irreducible component $\mathcal H_h$.
\item $p > p_0(H)$ so all degrees, separability, and local singularity types
are as in characteristic zero.
\end{itemize}

\begin{remark}[The $h=1$ exception]
For $h=1$ one has $(Y+1)^3 = (X+1)^3$, which splits into three lines.
Thus $h=1$ is handled separately and not included in the census clause.
\end{remark}

\subsection*{Genus bounds}

A curve of bidegree $(a,b)$ in $\mathbf P^1 \times \mathbf P^1$ has arithmetic genus $(a-1)(b-1)$.
Here $(a,b) = (3h,3k)$, so
\[
g_{\mathrm{arith}} = (3h-1)(3k-1).
\]

There are $hk$ pole-pole points $(X=-i, Y=-j)$.
A local expansion shows that each is an ordinary triple point,
hence contributes delta-invariant $3$.
Subtracting these gives
\[
g \le (3h-1)(3k-1) - 3hk \le 6hk.
\]
Similarly for $\mathcal H_h$ one obtains $g \le 6h^2$.

\subsection*{Complete point count}

\begin{lemma}
Let $\mathcal C$ be a cross component or $\mathcal H_h$.
Then
\[
\#\mathcal C(\Fp)
\le p + 1 + 12hk \sqrt{p} + 36 (hk)^2.
\]
\end{lemma}

\begin{proof}
Weil gives $p+1+2g\sqrt p \le p+1+12hk\sqrt p$.
Singular points are bounded by Bezout by $< 36(hk)^2$.
\end{proof}

\subsection*{Pairwise bounds}

For $h \ne k$,
\[
R_{h,k}
\le \frac{L_h L_k}{p}
+ 12hk \sqrt p
+ 36 (hk)^2
+ 4hk.
\]

For $2 \le h \le H$,
\[
R_{h,h}
\le \frac{L_h^2}{p}
+ L_h
+ 12h^2 \sqrt p
+ 36 h^4
+ 5h.
\]

For $h=1$,
\[
R_{1,1} \le 3 L_1 \le \frac{L_1^2}{p} + L_1 + p.
\]

\subsection*{Summation}

Let
\[
A_1 = \sum_{h\le H} h = \frac{H(H+1)}{2},
\quad
A_2 = \sum_{h\le H} h^2 = \frac{H(H+1)(2H+1)}{6}.
\]

Summing gives
\[
\Ncoinc(H)
\le \frac{S^2}{p} + S + p
+ 12 A_1^2 \sqrt p
+ 36 A_2^2
+ 4 A_1^2
+ 5 A_1.
\]

Using $A_1 \le H^2$, $A_2 \le H^3$,
\[
\Ncoinc(H)
\le \frac{S^2}{p} + S + p
+ 12 H^4 \sqrt p
+ 36 H^6
+ 4 H^4
+ 5 H^2.
\]

\subsection*{Main theorem}

\begin{theorem}
Assume CENSUS$_H$.
For
\[
H \le c_0 \, p^{1/6}, \quad c_0 = 1/4, \quad p \ge 4096,
\]
one has
\[
\boxed{
\Ncoinc(H) \le \frac{S^2}{p} + 4S.
}
\]
\end{theorem}

\begin{proof}
Since $S \ge pH/2$, we bound
\[
\frac{E}{S}
\le 24 \frac{H^3}{\sqrt p}
+ 72 \frac{H^5}{p}
+ 8 \frac{H^3}{p}
+ 10 \frac{H}{p}.
\]

For $H \le p^{1/6}/4$ this is $< 1/2$.
Also $p \le 2S$.
Thus
\[
\Ncoinc(H) \le \frac{S^2}{p} + 4S.
\]
\end{proof}

\subsection*{Additive-completion audit (with logs)}

Using additive completion,
\[
\sum_a |\widehat{1_I}(a)| \le p(2+\log p),
\]
and Bombieri bounds,
\[
|\sum e_p(aX+bY)| \le 40 hk \sqrt p.
\]

This gives per-pair
\[
R_{h,k}
\le \frac{L_h L_k}{p}
+ 64 hk \sqrt p (2+\log p)^2
+ 72 (hk)^2.
\]

Summing yields error
\[
E \sim H^4 \sqrt p (2+\log p)^2.
\]

Comparing with $S \sim pH$, we require
\[
H^3 (2+\log p)^2 / \sqrt p \ll 1.
\]

Thus the honest completion range is
\[
\boxed{
H \le c \, p^{1/6} / (2+\log p)^{2/3}.
}
\]

\begin{remark}
The log-loss is intrinsic to additive completion.
It cannot be removed by adjusting constants.
The full $p^{1/6}$ range holds only because
the windows are co-short and can be enlarged directly.
\end{remark}

\subsection*{Corollary}

\begin{corollary}
Under the theorem,
\[
\sum_{h \le H} C_h
\le \frac{S}{p} + 2\sqrt{S}
\le H + 2 \sqrt{pH}.
\]
Hence
\[
\sum_{h \le H} C_h
\le H + \sqrt{(4+o(1)) pH}.
\]
\end{corollary}

\subsection*{Limitation}

\begin{remark}
The trivial bound gives
\[
\sum_{h\le H} C_h \le \frac{3}{2} H^2.
\]
This dominates until $H \sim p^{1/3}$.
Since the theorem only holds up to $p^{1/6}$,
it does not improve the exponent.
Its value is structural: it certifies the
$S^2/p + O(S)$ coincidence shape.
\end{remark}

\subsection*{Assumptions beyond CENSUS$_H$}

\begin{enumerate}
\item Casoratian identity linking $\Delta_{r,h}$ and $\delta_h(r)$.
\item Degree data $\deg N_h = 3(h-1)$.
\item Genus bounds $g \le 6hk$ from bidegree and triple-point analysis.
\item Good reduction threshold $p_0(H)$.
\item Weil bound for curves over $\Fp$.
\item Bezout bound on singularities.
\item (Audit only) Bombieri additive bound and interval Fourier bound.
\item Explicit handling of $h=1$ outside the census.
\end{enumerate}

\end{document}
